\section{Einleitung}
\label{sec:introduction}

Wearable-Geräte sind aus dem Alltag nicht mehr wegzudenken. Smartwatches werden durchgehend am Handgelenk getragen und zeichnen kontinuierlich Bewegungsdaten über integrierte Inertialmesseinheiten (IMUs) auf, bestehend aus Beschleunigungssensoren und Gyroskopen. Diese Sensoren wurden ursprünglich für Schrittzählung und Fitness-Tracking entwickelt, erfassen jedoch ein weit breiteres Spektrum an Aktivitätsinformationen. Eine bislang kaum untersuchte Anwendung ist die automatische Erkennung von Schreibaktivität: die Frage, ob eine Person gerade mit der Hand schreibt, allein auf Basis der Handgelenksbewegung.

Handschreiben ist eine grundlegende Tätigkeit in Bildung, Beruf und Alltag. Im Vergleich zum Tippen auf einer Tastatur ist Handschreiben eine kognitiv anspruchsvollere motorische Aufgabe, die das Gehirn auf andere Weise beansprucht und nachweislich die Informationsverarbeitung fördert. Die Fähigkeit, Schreibaktivität passiv und unauffällig zu erkennen, ohne dass der Nutzer aktiv einen Timer starten oder mit dem Gerät interagieren muss, eröffnet vielfältige praktische Anwendungen. Dazu zählen die Erfassung produktiver Schreibzeit, das Monitoring kognitiver Aktivität in Bildungskontexten sowie der Aufbau smarter Umgebungen, die sich daran orientieren, ob eine Person gerade gedanklich etwas zu Papier bringt.

Bisherige Arbeiten zur Schreiberkennung mit handgelenksgetragenen Sensoren fokussieren sich vorwiegend darauf, \textit{was} geschrieben wird, also auf die Erkennung von Buchstaben, Wörtern oder Schreibgesten, anstatt der einfacheren, aber praktisch relevanteren Frage nachzugehen, \textit{ob} überhaupt geschrieben wird \cite{soubam2023, xiamotionhacker2018}. Diese Ansätze basieren häufig auf kontrollierten Laborbedingungen, wenigen Probanden und Evaluationsstrategien, die keine belastbare Aussage über die Generalisierbarkeit auf unbekannte Nutzer ermöglichen.

Die vorliegende Arbeit adressiert diese Einschränkungen. Wir erheben IMU-Daten einer Apple Watch Series~7 von 22 Probanden in einem ausbalancierten Studienprotokoll, das sowohl echte Schreibbedingungen als auch eine Reihe schwieriger Negativbeispiele umfasst, darunter Tastatur- und Smartphone-Tippen, Scrollen, Stiftfummeln und Gestikulieren. Die Ground-Truth-Labels werden ausschließlich während der Datenerhebung von einem Moleskine Smart Pen geliefert; nach dem Training läuft die Inferenz allein auf der Apple Watch. Alle Modelle werden mit Leave-One-Subject-Out-Kreuzvalidierung (LOSO) bewertet, bei der in jeder Falte ein Proband vollständig ausgelassen wird. Dies entspricht exakt der Generalisierungsanforderung eines realen Einsatzszenarios.

Unsere zentrale Forschungsfrage lautet:

\begin{quote}
\textit{Lässt sich Schreibaktivität allein aus den IMU-Daten einer Apple Watch zuverlässig erkennen, unabhängig von der schreibenden Person und dem geschriebenen Inhalt?}
\end{quote}

Wir vergleichen verschiedene Ansätze, darunter einen Random-Forest-Basisklassifikator mit 92 handkonstruierten Features, einen kausalen Hidden-Markov-Modell-Nachfilter (HMM) sowie tiefe Zeitreihenmodelle (TCN, GRU). Die experimentellen Ergebnisse und eine detaillierte Analyse der Modellleistung werden in Abschnitt~4 vorgestellt.

Der Rest des Artikels gliedert sich wie folgt: Abschnitt~2 bespricht verwandte Arbeiten zur Schreiberkennung mit handgelenksgetragenen Sensoren. Abschnitt~3 beschreibt das Datenerhebungsprotokoll, den Sensor-Setup und die Vorverarbeitungs-Pipeline. Abschnitt~4 präsentiert die experimentellen Ergebnisse und deren Analyse. Abschnitt~5 schließt mit einer Diskussion der Limitationen und zukünftiger Forschungsrichtungen.
